\subsectionaddtolist{الف) چسبندگی (Cohesion)}

چسبندگی به میزان ارتباط و انسجام اجزای داخلی یک ماژول یا کلاس اشاره دارد. به عبارت دیگر، نشان می‌دهد که عناصر یک واحد نرم‌افزاری تا چه حد در راستای یک هدف مشخص و مشترک عمل می‌کنند.

در طراحی مناسب، هر کلاس یا ماژول باید یک مسئولیت مشخص و متمرکز داشته باشد. در این حالت، گفته می‌شود چسبندگی بالا است که یک ویژگی مطلوب محسوب می‌شود. در مقابل، اگر یک کلاس چندین وظیفه نامرتبط را بر عهده داشته باشد، چسبندگی آن پایین بوده و این موضوع باعث کاهش خوانایی، نگهداری‌پذیری و قابلیت توسعه سیستم می‌شود.

\textbf{جمع‌بندی:}
چسبندگی بالا به معنای تمرکز هر واحد نرم‌افزاری بر یک هدف مشخص و واحد است.


\subsectionaddtolist{ب) اصل LSP در SOLID}

اصل جایگزینی لیسکوف بیان می‌کند که اشیاء یک کلاس پایه باید بتوانند بدون ایجاد اختلال در درستی برنامه، با اشیاء کلاس‌های مشتق‌شده جایگزین شوند.

به بیان دقیق‌تر، کلاس‌های فرزند باید رفتارهایی سازگار با قرارداد تعریف‌شده در کلاس پایه ارائه دهند و نباید پیش‌شرط‌ها را سخت‌تر یا پس‌شرط‌ها را ضعیف‌تر کنند. در غیر این صورت، استفاده از \lr{Polymorphism} منجر به بروز خطاهای منطقی در سیستم خواهد شد.

مثال معروف: فرض کنید کلاس \lr{Rectangle} را داریم. اگر بخوایم کلاس \lr{Square} را بسازیم و از \lr{Rectangle} ارث ببرد، به مشکل می‌خوریم! چون اگر عرض را تغییر دهیم، طول هم باید تغییر کند. اما \lr{Rectangle} انتظار دارد که ما می‌توانیم عرض و طول را جداگانه عوض کنیم. پس اگر جایی از کد \lr{Rectangle} استفاده می‌کنیم و ما به جای آن \lr{Square} بزاریم، منطق و کد به مشکل می‌خورد.

